\section{Execution labels, time limits and the generated test suite}\label{sup:labeling}

\subsection{Reading the query of an output}

Evaluation and execution
labeling read the SQL query of an output of the model in the same way.
If the output contains a fenced code block, its body is taken. The text
is then cut at the first semicolon, and what precedes the first
\texttt{SELECT} is dropped, except that a query beginning with
\texttt{WITH} keeps its \texttt{WITH} clause. An output without
\texttt{SELECT} is executed as it is, and a sample of which nothing
remains contains no query and is labeled incorrect. Every database is
opened read-only, so a statement that writes fails.

\subsection{Time limits}

Queries run under the time limit of the
benchmark's evaluation. On BIRD, one budget of \setting{timeout.bird}
covers a prediction followed by the gold SQL query, which is executed
with each prediction, and on Spider each query has
\setting{timeout.spider}. A query counts as reaching the limit only when
it used the processor for at least four fifths of the elapsed time. One
that reached it while waiting for the processor is run once more with no
other query running, and the outcome of that run counts. At
\setting{rounds.total} rounds, \fact{timeouts.bird.fedavg.4} predictions
of FedAvg SFT and \fact{timeouts.bird.k2.4} of \method{} reach the limit
on BIRD dev, and none on Spider test.

\subsection{Execution labels}

Execution labeling
(Section~\ref{sec:labeling}) applies the same limits and the same rule
to the samples of the training questions. A sample that exhausts its
limit while it runs is an execution error and is labeled incorrect. A
sample has no label when its result is never compared with the gold
result. A sample cut at the token limit is not executed
(\fact{train.bird.truncated.pref1} of the
\fact{train.bird.samples.pref1} samples of the first preference round on
BIRD). On BIRD, the time of the gold SQL query, measured once, counts
toward the budget of every sample of the question, and a completed
sample whose budget is thereby exceeded has no label. A gold SQL query
that takes more than half of the budget exceeds the limit for its
question. Its time is elapsed time, which depends on the load of the
machine, so the number of questions that the limit leaves without labels
differs between runs of the labeling step. On Spider, each query has its
own limit, and a completed sample is always compared with the gold
result.

\subsection{Questions that yield no pair}

A question whose gold SQL query
fails to execute or exceeds the time limit has no labeled sample. The
gold SQL query fails to execute for \fact{train.bird.gold_error}
questions of BIRD train and \fact{train.spider.gold_error} of Spider
train. On BIRD, it exceeds the time limit for
\fact{train.bird.gold_timeout.pref1} questions in the first preference
round of \method{} and \fact{train.bird.gold_timeout.pref2} in the
second, and on Spider for none. Spider train has
\fact{train.spider.no_rows.databases} databases without rows, which hold
\fact{train.spider.no_rows} questions, and BIRD train has none. On a
database without rows, queries of different meaning return the same
result, so its questions yield no pair. No evaluation set contains such
a database. All of these questions remain training targets in the
warm-up rounds.

\subsection{The generated test suite of BIRD dev}

BIRD has no official
test suite, and we generate one for BIRD dev with the fuzzing procedure
of the test-suite evaluation (Section~\ref{sec:datasets} of the main
paper cites it), which we reimplement: \fact{suite.variants} variants of
each of its \fact{suite.bird.databases} databases. A variant keeps the
schema of its database, and its tables are filled in foreign-key order
with random values of the column types, mixed with values of the source
database and with the literals of the gold SQL queries. Without a
distillation step, every variant is kept. A variant counts for a
question when the gold SQL query executes on it within the time limit,
returns at least one row, and has no result made ambiguous by ties under
\texttt{ORDER BY} with \texttt{LIMIT}. A prediction is correct under TS
when it is correct under EX on the original database and returns the
gold result on every variant that counts. Of the \fact{questions.bird}
questions of BIRD dev, \fact{suite.bird.counted} have at least one
variant that counts, and the other \fact{suite.bird.none} keep their EX
verdict.

\subsection{Agreement with an official suite}

On Spider dev, which has an
official suite, we generated a suite in the same way from its gold SQL
queries, with the same number of variants, and compared the verdicts of
the two suites on the predictions of our models on Spider dev
(\fact{suite.calib.files} sets of predictions). Of these predictions,
\fact{suite.calib.ex_correct} are correct under EX, and the two suites
agree on \fact{suite.calib.agreement_ex_correct}\% of them: the generated
suite rejects \fact{suite.calib.recall}\% of the
\fact{suite.calib.official_rejected} that the official suite rejects,
and \fact{suite.calib.ours_only} of the
\fact{suite.calib.official_accepted} that it accepts. The number of
variants was chosen among \fact{suite.calib.candidates} as the one with
the highest agreement on these predictions, so the agreement is measured
on the data that set it. No official suite exists against which the
suite of BIRD dev itself could be compared.
